\section{Conclusiones}

\begin{enumerate}
\item Los cinco algoritmos quedaron integrados en la clase \texttt{Clasificacion} del
framework, con un protocolo reproducible: preparación aprendida sólo con
entrenamiento, validación cruzada estratificada de 5 pliegues para elegir y una única
evaluación en prueba. Cada experimento se repite con un comando y deja toda la evidencia.

\item En Potabilidad, el mejor modelo es Random Forest con 500 árboles (F1 macro de
validación $0.644\pm0.022$; en prueba, F1 macro 0.659, exactitud 70.2\,\% y AUC 0.727).
Supera con claridad la referencia (59.6\,\% de exactitud) y a las demás familias, pero
detecta sólo el 43\,\% de las muestras potables. Las variables apenas difieren en
promedio entre clases y la señal aparece en los extremos y en la combinación de sulfato
y \texttt{ph}, algo que los ensambles de árboles profundos captan mejor que KNN o
AdaBoost con stumps.

\item En Diabetes, DT, RF, XGBoost y AdaBoost obtienen resultados equivalentes dentro de
la variación de la validación cruzada (F1 macro 0.80--0.84). RF fue el elegido por
validación; en prueba los cuatro detectan 8 de 15 diabéticos con AUC entre 0.89 y 0.93.
La glucosa estabilizada explica casi todo el poder predictivo, y los siete diabéticos
que ningún modelo detecta tienen glucosa normal en la medición, lo que marca un límite
de la información disponible.

\item Las variantes que mejoraron dependen del algoritmo y del dataset. Exigir hojas de al menos
10 observaciones mejoró el árbol de decisión en ambos datasets, pero limitar las hojas
de Random Forest empeoró el bosque, porque el promedio ya controla la varianza. Ponderar
las clases ayudó a los métodos de boosting cuando tendían a ignorar la clase minoritaria
(Potabilidad). Seleccionar atributos por información mutua ayudó
cuando la señal estaba en pocas variables con relación monótona (Diabetes) y perjudicó
cuando dependía de interacciones (Potabilidad).

\item La exactitud sola es engañosa con clases desbalanceadas. El F1 macro, el recall
por clase y la comparación con una referencia trivial fueron necesarios para interpretar
los resultados.
\end{enumerate}

Como trabajo futuro se propone elegir el umbral de decisión con validación cruzada para
aumentar el recall de la clase minoritaria, explorar combinaciones de hiperparámetros con
búsqueda en rejilla dentro de la validación cruzada y repetir el experimento con varias
semillas para estimar la estabilidad de las conclusiones.
